% CO1b (v1) -- Calculus of motion, PHYS 201, Fall 2024.
%
% Converted from "CO1b - Calculus Motion - v1 - Fall 2024 - Physics 201.docx",
% which is kept alongside this file as the source of record for the wording.
%
% This is the BLANK STUDENT version.  Its companion,
% CO1b-Fall2024-PHYS201-v1-solutions.tex, generates the same quiz with the
% worked solutions printed and records the answer key through randassign.
% The two files share their question code verbatim; keep them in step.
%
% To print a matching pair, set SEED to the same integer in both files (see the
% note in the pycode block below), then build each with the usual three passes:
%     pdflatex -interaction=nonstopmode <file>.tex
%     pythontex <file>.tex
%     pdflatex -interaction=nonstopmode <file>.tex

\documentclass[12pt]{exam}

\usepackage[margin=1in]{geometry}

\newcommand{\class}{Physics 201}
\newcommand{\term}{Fall 2024}
\newcommand{\examnum}{CO1b (v1)}
\newcommand{\examdate}{\today}
\newcommand{\timelimit}{50 Minutes}


% Engine-specific settings
% Detect pdftex/xetex/luatex, and load appropriate font packages.
% This is inspired by the approach in the iftex package.
% pdftex:
\ifx\pdfmatch\undefined
\else
    \usepackage[T1]{fontenc}
    \usepackage[utf8]{inputenc}
\fi
% xetex:
\ifx\XeTeXinterchartoks\undefined
\else
    \usepackage{fontspec}
    \defaultfontfeatures{Ligatures=TeX}
\fi
% luatex:
\ifx\directlua\undefined
\else
    \usepackage{fontspec}
\fi
% End engine-specific settings

\usepackage{amsmath,amssymb}
\usepackage{fullpage}
\usepackage{graphicx}
\usepackage[svgnames]{xcolor}
\usepackage{url}
\urlstyle{same}

\usepackage{pythontex}
% \restartpythontexsession{\thesection}


\usepackage[framemethod=TikZ]{mdframed}

\newcommand{\pytex}{Python\TeX}
\renewcommand*{\thefootnote}{\fnsymbol{footnote}}

\usepackage{siunitx}
\usepackage{pgfplots}

% The Word original carries the class/quiz line with the name and id# rules as a
% running header, and the work note as a running footer.  The exam class does
% both directly, so `nopageno` is not loaded here: this page style already
% suppresses the page number.
\pagestyle{headandfoot}
\header{\textbf{\class} -- \textbf{\examnum}}{}%
       {\textbf{Name:}~\makebox[2in]{\hrulefill}~\textbf{id\#}~\makebox[1.1in]{\hrulefill}}
\headrule
\footer{}{\textbf{SHOW ALL WORK FOR HIGHEST RATING}}{}

% The exam class puts the head rule close enough to the text that a question
% starting at the top of a page touches it.
\setlength{\headsep}{0.35in}


\begin{document}

% The quiz body is generated from the single pycode block below, so every number
% on the page is computed rather than typed.  Question 1 is deliberately fixed;
% question 2 is drawn fresh for each paper.

\begin{pycode}
import random

import sympy as sp

# ---------------------------------------------------------------------------
# Which paper this build produces
# ---------------------------------------------------------------------------
# Leave SEED as None for a fresh paper on every build.  Set it to an integer --
# the SAME integer here and in the solutions file -- to reproduce one exact
# paper, so the blank quiz and the worked key describe the same particle.
#
# Only the standard-library RNG is used, so seeding controls every draw in this
# document.  (`from pylab import *` would shadow `uniform` with the numpy
# generator, which random.seed() does not drive; see CLAUDE.md.)
SEED = None

paper_code = SEED if SEED is not None else random.randrange(1, 10**6)
random.seed(paper_code)

t = sp.Symbol('t', nonnegative=True)
tau = sp.Symbol('tau', nonnegative=True)


# ---------------------------------------------------------------------------
# Question 1 -- Artemis rocket.  Deliberately NOT randomised.
# ---------------------------------------------------------------------------
V_BURNOUT = 1800        # m/s reached at the end of the burn
T_BURN_MINUTES = 2      # "the first two minutes after liftoff"
MINUTE_WORDS = {1: 'one', 2: 'two', 3: 'three', 4: 'four', 5: 'five'}


def question_rocket():
    """Average acceleration over the burn, then the height reached.

    The rocket starts from rest and the acceleration is taken as constant, so
    v(t) = a t and the height is the integral of v across the burn.  Both are
    evaluated with sympy rather than by quoting a kinematic formula.
    """
    burn = T_BURN_MINUTES * 60                        # seconds
    a_avg = sp.Rational(V_BURNOUT, burn)              # (v - v0)/(t - t0), v0 = 0
    v_of_t = sp.integrate(a_avg, (tau, 0, t))         # = a t
    height = sp.integrate(v_of_t, (t, 0, burn))       # = a t^2 / 2

    word = MINUTE_WORDS[T_BURN_MINUTES]
    statement = (
        r'The Artemis rocket is estimated to be able to reach '
        r'$\SI{%.4g}{\meter\per\second}$ in the first %s minutes after liftoff.'
    ) % (V_BURNOUT, word)
    parts = [
        r'What is the average acceleration of the rocket during this time?',
        r'If air resistance were not a factor and if the rocket went straight '
        r'up, how high off the ground would it be at the end of the first %s '
        r'minutes of its flight?' % word,
    ]
    return statement, parts, (a_avg, v_of_t, height, burn)


# ---------------------------------------------------------------------------
# Question 2 -- particle on the x axis.  Randomised.
# ---------------------------------------------------------------------------
# x(t) = A t - B t^3, so v = A - 3B t^2 vanishes at t = sqrt(A/(3B)).  Drawing B
# and that instant, then setting A = 3 B t0^2, keeps the root a whole number of
# seconds the way the original paper's 24t - 2.0t^3 does.  The original draw
# (A = 24, B = 2, t0 = 2) is one of the possibilities.
def question_particle():
    b_halves = random.choice([2, 3, 4, 5, 6])         # B = 1.0 .. 3.0
    coeff_b = sp.Rational(b_halves, 2)
    t_zero = random.choice([2, 3])                    # instant when v = 0
    coeff_a = 3 * coeff_b * t_zero**2

    # Evaluate part (b) somewhere other than the instant found in part (a), so
    # the two parts cannot collapse onto the same answer.
    while True:
        t_eval = sp.Rational(random.choice([2, 3, 4, 5, 6]), 2)
        if t_eval != t_zero:
            break

    x_of_t = coeff_a * t - coeff_b * t**3
    v_of_t = sp.diff(x_of_t, t)
    a_of_t = sp.diff(v_of_t, t)

    roots = [r for r in sp.solve(sp.Eq(v_of_t, 0), t) if r.is_real and r >= 0]
    t_rest = max(roots)
    a_at_rest = a_of_t.subs(t, t_rest)
    a_at_eval = a_of_t.subs(t, t_eval)

    statement = (
        r'A particle moving along the $x$ axis has a position given by '
        r'$x = (%.4g\,t - %.1f\,t^{3})$ m, where $t$ is measured in s.'
    ) % (float(coeff_a), float(coeff_b))
    parts = [
        r'What is the magnitude of the acceleration of the particle at the '
        r'instant when its velocity is zero?',
        r'What is the acceleration of the particle at '
        r'$t = \SI{%.1f}{\second}$?' % float(t_eval),
    ]
    working = (x_of_t, v_of_t, a_of_t, t_rest, abs(a_at_rest), t_eval, a_at_eval)
    return statement, parts, working


# ---------------------------------------------------------------------------
# Layout
# ---------------------------------------------------------------------------
WORKSPACE = r'\vspace{2.1in}'


def emit(statement, parts):
    """One numbered question with its lettered parts and room to work."""
    print(r'\question ' + statement)
    print(r'\begin{parts}')
    for part in parts:
        print(r'\part ' + part)
        print(WORKSPACE)
    print(r'\end{parts}')


rocket_statement, rocket_parts, _ = question_rocket()
particle_statement, particle_parts, _ = question_particle()

# The paper code identifies which draw this copy came from; setting SEED to it
# in the solutions file reprints the matching key.
print(r'\hfill{\scriptsize Paper code: %d}\par' % paper_code)
print(r'\setcounter{question}{0}')
print(r'\begin{questions}')
emit(rocket_statement, rocket_parts)
print(r'\newpage')
emit(particle_statement, particle_parts)
print(r'\end{questions}')
\end{pycode}

\end{document}
